% !TeX root = ../../main.tex
\subsubsection{simCRISPR interaction benchmark}

For the factorial interaction benchmark, we used simCRISPR
\citep{zhu2026simcrispr}, which simulates knockout induction crossed with
treatment, propagated over several days of cell growth, with PCR amplification
and sequencing applied afterwards.  Each sgRNA has a knockout effect, a
treatment effect, and an interaction estimated with an explicit design-matrix
column.

Each of three seeds generated 2,000 sgRNAs, of which 300 were non-targeting
and 200 were safe-harbor controls, in three independent replicates of each of
the four arms, so twelve libraries and eight residual degrees of freedom.
Growth was logistic.  The alternative exponential model is unbounded: over
five days a single guide reached a fifth of the library and library sizes
spanned a 175-fold range, which does not represent a pooled screen.  Initial
library columns were dropped, since they are not part of the factorial
contrast.  Amplification used the package defaults and each library was
sequenced to 500 reads per guide.  BARCS fitted
$\logit(\mu_{gi})=\beta_0+\beta_1 k_i+\beta_2 t_i+\beta_3 k_i t_i$ for
knockout indicator $k_i$ and treatment indicator $t_i$, and tested $\beta_3$.

Three denominators were compared.  The library denominator is the full column
total, BARCS's default.  The two control denominators hold a chosen control
class at a constant share, rescaled to the mean library size and rounded to
produce effective library totals.  These normalized totals define a
control-relative analysis.  Safe-harbor guides are cut and carry the
DNA-damage response; non-targeting guides undergo neither.  We split each
control class in half, using one half for normalization and calibration and
holding out the other for evaluation.  No guide used to fit the null was
evaluated against that fit.  MAGeCK-MLE received the same design matrix with
an explicit interaction column, the same calibration-half controls, and one
sgRNA per gene, because the simulated truth is guide-level.

CRISPhieRmix was not run on these data.  It borrows strength across the guides
within a gene, whereas simCRISPR assigns every sgRNA an independent
interaction, so grouping guides into synthetic genes would impose the
structure that the model exploits.

Every targeting guide receives a nonzero interaction, most too small to
resolve.  We therefore fixed a magnitude before scoring: guides with
$\lvert\mathrm{INT}\rvert\geq0.2$ are the positives and the held-out
non-targeting guides, whose interaction is exactly zero, are the negatives.
Targeting guides below that magnitude are scored for effect recovery but
excluded from the detection metrics.  Effect recovery is the Spearman
correlation between the estimated interaction coefficient and the simulated
one across all targeting guides, which needs no threshold.  The held-out
control call rate and the safe-harbor call rate are reported separately: the
first measures false positives against a ground-truth zero, the second
measures how often the cutting response alone is mistaken for an interaction.
